\begin{examplebox}
\textbf{例1：右开口实际切点弦。}
对 $y^2=4x$，$P=(-1,0)$ 满足 $y_0^2-4x_0=4>0$。
切点弦 $0=2(x-1)$，即 $x=1$，切点为 $(1,\pm2)$。
两条切线 $y=x+1$、$y=-x-1$ 都经过 $P$。
\medskip
\textbf{例2：原题给定点不能引两条实切线。}
对 $x^2=6y$，$P=(-1,1)$，$x_0^2-6y_0=-5<0$，故没有实切点弦。
形式极线为 $x+3y+3=0$；与抛物线联立得
$x^2+2x+6=0$，判别式 $4-24=-20<0$。
\medskip
\textbf{例3：由切点弦反求极点。}
对 $y^2=4x$，若切点弦为 $y=x-1$，比较 $y_0y=2(x+x_0)$，
得 $P=(-1,2)$。检验 $y_0^2-4x_0=8>0$，所以确有两条实切线。
\end{examplebox}
